% Crust benchmark report.
%
% Static shell. The data-driven bodies are regenerated by the plot scripts
% from the latest results/*.json:
%
%   crust.tex             hand-written intro (tools/crust.tex)
%   crust_bench_body.tex  benchmarks/plot_crust.py    -- Crust vs gcc
%   report_body.tex       benchmarks/plot_rpython.py  -- cross-runtime
%   minipy_body.tex       benchmarks/plot_minipy.py   -- minipy interpreter
%
% Build with `make benchmarks_report`, which runs each harness, renders the
% figures, and invokes pdflatex in the figure output directory.
\documentclass[11pt]{article}
\usepackage[margin=1in]{geometry}
\usepackage{graphicx}
\usepackage{float}
\usepackage{booktabs}
\usepackage{array}
\usepackage{xcolor}
\usepackage{hyperref}
\usepackage{orcidlink}

\hypersetup{colorlinks=true, linkcolor=blue, urlcolor=blue}

\title{Crust: The C + C++ Rust Multi-Compiler}
\author{B.S.Hartshorn \orcidlink{0009-0004-2853-655X} \small\url{https://github.com/brentharts/crust}}
\date{\today}

\begin{document}
\maketitle
\tableofcontents
\newpage

\section{What Crust is}

Crust is a Rust subset that compiles through a C compiler. A single
translation unit may hold C, Rust and rpython at once, and after lowering
there is no boundary between them: everything becomes C, and the C compiler
sees one program.

\begin{verbatim}
#include "schemes.py"        /* rpython: string and list work         */
int printf(const char *, ...);
struct Frames { words: [u64; 8], free: i64 }   /* Rust: fixed layouts  */
impl Frames {
    fn alloc(&mut self) -> i64 { ... }
}
int main(void) { ... }                          /* C: entry point      */
\end{verbatim}

The point of mixing three languages is that each is short at something
different. rpython has lists, strings and a `\%'-formatting runtime, so URL
parsing and table building are a few lines there and a rewrite in the Rust
subset. Rust has fixed layouts, traits and generics with no allocator
underneath, so a bitmap allocator or a scheduler is natural. C is the entry
point and the linkage surface.

Nothing is copied at the boundaries. \texttt{tools/py2c.py} lowers a typed
rpython list to three words --- \texttt{\{ int* data; long len; long cap; \}}
--- and \texttt{PyList<T>} in Crust's bundled core has exactly that layout, so
a list rpython built is walked directly from Rust with a pointer cast and no
conversion.

\subsection{What is in the subset}

Structs, enums (including data-carrying ones, lowered to tagged unions),
traits with default methods and associated consts and types, generics by
monomorphisation, closures lifted to static functions, slices, \texttt{match}
with bindings, \texttt{?}, \texttt{impl} blocks, operator sugar for the
bundled containers, and a small core: \texttt{Vec<T>}, \texttt{Box<T>},
\texttt{String}, \texttt{Formatter} with \texttt{write!} and \texttt{format!},
\texttt{Rc}/\texttt{Arc}, \texttt{VecDeque<T>}, and the \texttt{core} free
functions real code reaches for (\texttt{slice::from\_raw\_parts},
\texttt{ptr::null\_mut}, \texttt{cmp::min}).

Generics are monomorphised, so a trait call is a direct call and an
instantiation costs what a hand-written version would --- the
\texttt{traits} and \texttt{generics} benchmarks below are there to check that
claim rather than assert it.

\subsection{What is deliberately absent}

There is no borrow checker, no \texttt{dyn Trait}, no iterator protocol, and
no \texttt{Drop}. Each of those is a decision rather than an omission waiting
to be filled in:

\begin{itemize}
  \item \textbf{No \texttt{Drop}} means every allocating type has an explicit
    \texttt{free\_buf}. Scope exit cannot run code, so pretending otherwise
    would leak silently.
  \item \textbf{No iterator protocol}, because rpython already has one and its
    output is a plain struct. Iteration-heavy code is written on that side.
  \item \textbf{\texttt{Mutex<T>} does not lock.} It exists so source written
    against \texttt{std} parses. Crust has no threads --- no spawn, no atomics,
    no memory model --- so there is nothing for a lock to protect against, and
    a lock that does nothing is consistent with the rest of the model rather
    than a hole in it. The moment real concurrency exists it becomes dangerous
    and must be replaced, not extended.
\end{itemize}

\section{CrustOS, and what it borrows from Redox}

CrustOS is a small operating system prototype assembled from the parts of
\href{https://redox-os.org}{Redox} that Crust can compile, plus the parts
supplied here. \texttt{tools/crustos.py} clones the Redox kernel and relibc,
compiles the compatible subset, and links it with a scheme layer written in
rpython and a scheduler written in Rust.

\begin{verbatim}
python3 tools/crustos.py fetch    # clone redox-os/kernel and relibc
python3 tools/crustos.py run      # compile the subset, link, run
\end{verbatim}

\subsection{Redox Compatibility}

Of 525 \texttt{.rs} files in the Redox kernel and relibc, Crust lowers items in
about 110 and around 94 of those compile to objects, of which 89 link
together. The upstream code that survives is mostly the memory manager and the
architecture definitions --- \texttt{rmm}'s page tables and flags, the buddy
and bump frame allocators, syscall numbers, layout constants. The
\texttt{heap offset} CrustOS prints at startup is upstream's
\texttt{kernel\_heap\_offset()}, called from code compiled here.

Five of the ninety-four objects cannot join the link, for two reasons that
both follow from compiling files never meant to be one program: Redox ships
one implementation per architecture and Crust flattens module paths, so the
\texttt{aarch64} and \texttt{riscv64} versions of a function collide; and some
files reference symbols from parts of Redox that do not compile.
\texttt{tools/crustos.py} selects a linkable subset with \texttt{nm} and
reports what it dropped and why.

\subsection{Redox Paging}

Redox \texttt{rmm} writes its page flags generically over the architecture, reading
constants off a trait:

\begin{verbatim}
impl<A: Arch> PageFlags<A> {
    pub fn new() -> Self {
        Self::from_data(A::ENTRY_FLAG_DEFAULT_PAGE | A::ENTRY_FLAG_NO_EXEC)
    }
}
\end{verbatim}

CrustOS uses the same shape, and Crust monomorphises it: one implementation of
the flag logic, one instantiation per architecture, every constant resolved at
compile time with no dispatch.



\section{Crust benchmarks}

These compare code generated by ShivyCX against gcc on \emph{identical C
input}: the front end lowers each source to C once, and that same C is then
compiled by both. What is measured is therefore code generation alone --- the
front end is shared and cancels out. Both front ends appear here: a
\texttt{.rs} benchmark is lowered by \texttt{shivyc.crust}, a \texttt{.cpp}
one by \texttt{tools.cpprust}, and past that point they are the same thing,
since both emit plain C over struct pointers.

Each benchmark isolates one part of the lowering rather than blending them, so
a regression in (say) tag dispatch shows up as one bar rather than being
averaged away. The benchmarks do not share a workload, so bar heights are
comparable \emph{within} a benchmark, across compilers --- not across
benchmarks.

\subsection*{What virtual dispatch costs}

\texttt{cpp\_methods} and \texttt{cpp\_dispatch} are a matched pair: the same
arithmetic, the same iteration count, the same runtime alternation between two
objects, the same call through a pointer. The only difference is that
\texttt{mix} is \texttt{virtual} in one and not the other, so the ratio
between them isolates dispatch from loop shape, cache behaviour and branch
prediction.

Two costs are folded into that ratio. One is the indirect call, which any
vtable scheme pays. The other is the thunk: a vtable slot's \texttt{this} is
the class that first declared the method, the implementation takes its own,
and a small forwarding function converts between them. Casting the function
pointer instead would remove that hop --- it is what a real C++ compiler
emits --- but it is undefined behaviour in C, and the premise of this lowering
is that its output is C. The gap between the two bars is the price of that
choice.

\IfFileExists{crust_bench.pdf}{%
  \begin{figure}[H]\centering
  \includegraphics[width=\linewidth]{crust_bench.pdf}
  \caption{Runtime (top) and binary size (bottom) for each Crust benchmark,
  built by ShivyCX and by gcc at three optimisation levels. ShivyCX emits
  \texttt{-O0}-class code, so the useful comparisons are against
  \texttt{gcc -O0}; the \texttt{-O2}/\texttt{-O3} bars show how much is left
  on the table.}
  \end{figure}}{%
  \emph{Crust figures not built --- run \texttt{python3
  benchmarks/run\_crust\_benchmarks.py} first.}}

\IfFileExists{crust_bench_body.tex}{\input{crust_bench_body.tex}}{}

\section{Runtime comparisons}

\IfFileExists{minipy_body.tex}{\input{minipy_body.tex}}{%
  \emph{MiniPy runtime figures not built --- run \texttt{python3
  benchmarks/run\_minipy\_benchmarks.py} first.}}

\section{What is measured}

Each benchmark in \texttt{benchmarks/rpython/} is a single \emph{pure-Python}
program. The same source is executed by four backends, which must all agree on
the program's exit code (a differential-correctness check):

\begin{itemize}
  \item \textbf{CPython} --- the program run directly under \texttt{python3}.
  \item \textbf{PyPy3} --- the same program under the PyPy3 JIT (omitted if PyPy3
        is not installed).
  \item \textbf{py2c$+$gcc} --- the program transpiled to C by
        \texttt{tools/py2c.py} and compiled with \texttt{gcc -O2}.
  \item \textbf{self-hosted} --- the \emph{same} transpiled C compiled by
        \texttt{shivyc\_native}, the ShivyCX compiler built from its own
        sources. This is the headline column: ShivyCX compiling the C that its
        own \texttt{.py}$\rightarrow$C translator produced.
\end{itemize}

For every backend three quantities are recorded:

\begin{itemize}
  \item \textbf{Runtime} --- best of five wall-clock measurements of the program
        proper.
  \item \textbf{Peak memory} --- the exact maximum resident set size of the
        process, read from the child's \texttt{rusage} (\texttt{ru\_maxrss}) by
        a $\sim$1.5\,MB C probe, so the harness's own footprint does not inflate
        it.
  \item \textbf{Compile time} --- for the two compiled backends, the wall-clock
        time the C compiler (gcc or \texttt{shivyc\_native}) spent on the py2c
        output. The interpreters have no compile step.
\end{itemize}

ShivyCX emits unoptimised, \texttt{-O0}-class code, so the self-hosted runtime is
expected to trail \texttt{gcc -O2}; the interesting questions are how it compares
to the \emph{interpreters} (it beats them on compute-bound, allocation-light
programs and trails them where boxed-list element access dominates) and how its
compile speed and the resulting binary's memory footprint compare to gcc.

\IfFileExists{report_body.tex}{\input{report_body.tex}}{%
  \emph{rpython cross-runtime figures not built. This section needs the
  self-hosted native compiler, which \texttt{tools/selfhost.py compiler}
  builds --- run \texttt{python3 benchmarks/run\_rpython\_benchmarks.py}
  first.}}

\section{Compiler features versus gcc}

These measure ShivyCX language extensions that gcc has no equivalent for, each
against the same program built with the feature off and with \texttt{gcc -O0}.
They are a different question from the runtime comparisons above: not ``how
fast is the generated code'', but ``what does the extension buy''.

\IfFileExists{benchmarks.png}{%
  \begin{figure}[H]\centering
  \includegraphics[width=\linewidth]{benchmarks.png}
  \caption{Feature benchmarks: ShivyCX with the feature off, on, and gcc -O0.}
  \end{figure}}{%
  \emph{Feature figures not built -- run \texttt{python3
  tools/crust\_benchmarks.py features} first.}}

\IfFileExists{benchmarks2.png}{%
  \begin{figure}[H]\centering
  \includegraphics[width=\linewidth]{benchmarks2.png}
  \caption{Memory-safety detection: which errors each compiler catches.}
  \end{figure}}{}

\section{Metamorphic returns}

A \texttt{\_\_metamorphic\_\_} function does not keep its return address on the
stack. The caller writes the address the callee should come back to, and the
callee ``returns'' with an indirect jump to it. The point is to avoid the
stack traffic and the return predictor entirely --- but done carelessly it is
far \emph{slower} than an ordinary call, so the interesting question is which
lowering is correct and when the feature actually pays off.

The right lowering is a register-indirect jump whose target has been patched
into a \texttt{mov} immediate, with that patch \emph{hoisted} out of any loop
whenever the call graph shows the return site is loop-invariant. Routing the
return through a memory slot instead pays a load; patching the immediate on
every call is a self-modifying-code machine clear and costs hundreds of cycles.

\IfFileExists{metamorphic_bench.pdf}{%
  \begin{figure}[H]\centering
  \includegraphics[width=\linewidth]{metamorphic_bench.pdf}
  \caption{Left: cycles per call for each return lowering (log axis). Right:
  call/ret versus hoisted metamorphic as the call chain deepens.}
  \end{figure}}{%
  \emph{Metamorphic figures not built --- run \texttt{python3
  benchmarks/run\_metamorphic\_benchmarks.py} first.}}

\IfFileExists{metamorphic_body.tex}{\input{metamorphic_body.tex}}{}

Unlike the classic way a compiler removes call overhead --- inlining --- this
does not grow with the callee. Inlining a large or deeply-nested function
explodes code size and instruction-cache pressure; the metamorphic return is a
few bytes regardless of callee size, and its win over \texttt{call}/\texttt{ret}
appears exactly where inlining and the return predictor both struggle: chains
deeper than the return-stack buffer. See \texttt{METAMORPHIC.md} for the full
discussion and the per-thread specialisation plan for multi-call-site
functions.

\section{Reproducing}

\begin{verbatim}
make benchmarks-report      # harness -> figures -> this PDF
\end{verbatim}

\noindent
The harness builds \texttt{shivyc\_native} once (or reuses one supplied via
\texttt{\$SHIVYC\_NATIVE}); results land in
\texttt{benchmarks/results/rpython\_results.json} and the figures in
\texttt{\$BENCH\_PLOT\_DIR} (default \texttt{/tmp/shivyc\_benchmarks}).


\begin{thebibliography}{99}

\bibitem{vixra2607}
B.S. Hartshorn, \emph{Beyond WebAssembly: Rethinking the Web Browser with
Post-JavaScript, VM-Free Page Execution}, 2026.
\url{https://ai.vixra.org/abs/2607.0013}

\bibitem{vixra2606}
B.S. Hartshorn, \emph{Rethinking Meta-Interpreters for High-Performance
Execution}, 2026. \url{https://ai.vixra.org/abs/2606.0084}

\bibitem{vixra2507}
B.S. Hartshorn, \emph{Foundational Problems with Compilers and Operating
Systems}, 7 July 2025. \url{https://ai.vixra.org/abs/2507.0081}

\end{thebibliography}

\end{document}
